% !TEX root = ../algebra_02.tex

\section{Группа автоморфизмов поля из $p^n$ элементов}

\begin{Theorem}{Группа автоморфизмов конечного поля}{}
	Группа автоморфизмов $\operatorname{Aut}(\mathbb{F}_{p^n})$ является циклической группой порядка $n$, порожденной автоморфизмом Фробениуса $\Phi(x) = x^p$.
\end{Theorem}

\begin{proof}
	Докажем в два этапа.
	\begin{enumerate}
		\item Покажем, что порядок $\Phi$ равен $n$.
		Для любого $a \in \mathbb{F}_{p^n}$ имеем $\Phi^n(a) = a^{p^n} = a$, откуда $\Phi^n = \operatorname{id}_{\mathbb{F}_{p^n}}$.
		Допустим, существует $m < n$ такое, что $\Phi^m = \operatorname{id}_{\mathbb{F}_{p^n}}$. Тогда для любого $a \in \mathbb{F}_{p^n}$ выполнено $a^{p^m} - a = 0$.
		Получается, что многочлен $X^{p^m} - X$ имеет $p^n$ различных корней. Это невозможно, так как его степень $p^m$ строго меньше $p^n$. Значит, порядок $\Phi$ в точности равен $n$.

		\item Докажем, что $|\operatorname{Aut}(\mathbb{F}_{p^n})| \le n$.
		Предположим противное: существуют различные автоморфизмы $\psi_1, \dots, \psi_{n+1} \in \operatorname{Aut}(\mathbb{F}_{p^n})$. Мультипликативная группа конечного поля циклическая, поэтому $\mathbb{F}_{p^n}^* = \langle a \rangle$ для некоторого элемента $a$.

		Пусть $f = \operatorname{Irr}_{\mathbb{F}_p}(a)$ --- его минимальный многочлен над простым подполем $\mathbb{F}_p$, причем $\deg f = n$. Заметим, что любой автоморфизм $\psi_j$ сохраняет элементы $\mathbb{F}_p$ неподвижными (так как $\psi_j(1)=1$).

		Поскольку $f(a) = 0$ и коэффициенты $f$ лежат в $\mathbb{F}_p$, то:
		\[ \forall j : \quad \psi_j(f(a)) = 0 \implies f(\psi_j(a)) = 0 \]
		Таким образом, элементы $\psi_1(a), \dots, \psi_{n+1}(a)$ являются корнями многочлена $f$. Так как степень $f$ равна $n$, по принципу Дирихле найдутся индексы $i \neq j$, для которых $\psi_i(a) = \psi_j(a)$.
		Но элемент $a$ порождает все поле $\mathbb{F}_{p^n}$, поэтому совпадение автоморфизмов на порождающем элементе влечет их полное совпадение: $\psi_i = \psi_j$. Мы получили противоречие.
	\end{enumerate}
	Следовательно, группа автоморфизмов состоит в точности из $n$ элементов $\{\operatorname{id}, \Phi, \Phi^2, \dots, \Phi^{n-1}\}$ и является циклической.
\end{proof}

\begin{Statement}{Соответствие Галуа для конечных полей}{}
	Существует взаимно однозначное соответствие (две взаимно обратные биекции) между подполями $F$ расширения (где $\mathbb{F}_p \subset F \subset \mathbb{F}_{p^n}$) и подгруппами $H \le \operatorname{Aut}(\mathbb{F}_{p^n})$:
	\begin{itemize}
		\item Каждому подполю $F$ ставится в соответствие подгруппа $H = \operatorname{Aut}(\mathbb{F}_{p^n}/F)$.
		\item Каждой подгруппе $H$ ставится в соответствие неподвижное подполе $L^H = \{ a \in \mathbb{F}_{p^n} \mid \forall h \in H : h(a) = a \}$.
	\end{itemize}
\end{Statement}
